% !TeX root = ../sections/02_exercises.tex

\begin{exercise}
	R-4. $R$ は可換環とする。
	$F$ が階数 $n$ の自由 $R$-加群であれば，
	\[
	\operatorname{Hom}_R(F,R)
	\]
	も階数 $n$ の自由 $R$-加群であることを示せ。
\end{exercise}

\begin{background}
	$F$ が階数 $n$ の自由 $R$-加群であるとは，ある基底
	\[
	e_1,\dots,e_n\in F
	\]
	が存在して，任意の $x\in F$ が一意的に
	\[
	x=a_1e_1+\cdots+a_ne_n
	\qquad
	(a_1,\dots,a_n\in R)
	\]
	と表されることをいう。
	
	$\operatorname{Hom}_R(F,R)$ は，$F$ から $R$ への
	$R$-加群準同型全体であり，これを $F$ の双対加群という。
	
	基底 $e_1,\dots,e_n$ に対して，双対基底
	\[
	e_1^\ast,\dots,e_n^\ast\in\operatorname{Hom}_R(F,R)
	\]
	を
	\[
	e_i^\ast(e_j)=
	\begin{cases}
		1 & (i=j),\\
		0 & (i\neq j)
	\end{cases}
	\]
	によって定める。
\end{background}

\begin{strategy}
	$F$ の基底を
	\[
	e_1,\dots,e_n
	\]
	とする。
	これに対して双対基底
	\[
	e_1^\ast,\dots,e_n^\ast
	\]
	を定義する。
	
	任意の $\varphi\in\operatorname{Hom}_R(F,R)$ に対して，
	$\varphi$ は基底の値
	\[
	\varphi(e_1),\dots,\varphi(e_n)
	\]
	によって完全に決まる。
	
	実際，任意の
	\[
	x=a_1e_1+\cdots+a_ne_n
	\]
	に対して，
	\[
	\varphi(x)
	=
	a_1\varphi(e_1)+\cdots+a_n\varphi(e_n)
	\]
	である。
	
	したがって，
	\[
	\varphi
	=
	\varphi(e_1)e_1^\ast+\cdots+\varphi(e_n)e_n^\ast
	\]
	と表せる。
	また，この表示は一意的である。
	
	よって，
	\[
	e_1^\ast,\dots,e_n^\ast
	\]
	は $\operatorname{Hom}_R(F,R)$ の基底である。
	したがって，
	\[
	\operatorname{Hom}_R(F,R)
	\]
	は階数 $n$ の自由 $R$-加群である。
\end{strategy}